\documentclass[12pt,leqno]{article}

\usepackage{amsmath}
\usepackage{amssymb}
\usepackage{enumitem}
\usepackage{hyperref}
\usepackage{logicproof}

\setlength{\textheight}{9.1in}
\setlength{\topmargin}{-.1in}
\setlength{\headsep}{0in}
\setlength{\oddsidemargin}{-.1in}
\setlength{\textwidth}{6.5in}

\setlength{\parskip}{1em}
\setlength{\parindent}{0pt}

% for logicproof
\setlength{\subproofhorizspace}{.5em}
\setlength{\intersubproofvertspace}{.5em}

% custom symbols
\newcommand{\DefEq}{\stackrel{\text{def}}{=}}

% proof rules
\newcommand{\Intro}[1]{{#1}{\text{I}}}
\newcommand{\Elim}[1]{{#1}{\text{E}}}
\newcommand{\ElimA}[1]{{#1}{\text{E}_1}}
\newcommand{\ElimB}[1]{{#1}{\text{E}_2}}

% truth tables
\newcommand{\T}{\mathtt{T}}
\newcommand{\F}{\mathtt{F}}

\title{CS 511 Homework \#2}
\author{Anthony DeRossi}
\date{19 Sep 2024}

\begin{document}
\maketitle

\section*{Exercise 1}

Go to page 9 in \textbf{\textit{Lecture Slides 06}}. Your task is to carefully
write all the details of the proof by \textit{structural induction}.

\medskip
Proposition: For all strings $s, t \in A^*$, it holds that $\text{reverse}(s
\cdot t) = \text{reverse}(t) \cdot \text{reverse}(s)$.

Proof: Use structural induction on $t \in A^*$ to prove the property $P(t)$
defined by:
\[
  P(t) \DefEq \text{``for all $s \in A^*$ it holds that
    $\text{reverse}(s \cdot t) = \text{reverse}(t) \cdot \text{reverse}(s)$''}
\]

\medskip
Base case: $t = \varepsilon$
\[
  P(\varepsilon) \DefEq \text{``for all $s \in A^*$ it holds
    that $\text{reverse}(s \cdot \varepsilon) = \text{reverse}(\varepsilon)
    \cdot \text{reverse}(s)$''}
\]

\begin{enumerate}
\item $\text{reverse}(s \cdot \varepsilon) = \text{reverse}(\varepsilon)
  \cdot \text{reverse}(s)$
\end{enumerate}

By $\text{reverse}(\varepsilon) \DefEq \varepsilon$,

\begin{enumerate}[resume]
\item $\text{reverse}(s \cdot \varepsilon) = \varepsilon \cdot
  \text{reverse}(s)$
\end{enumerate}

By the definition $\text{reverse}(s \cdot x) \DefEq x \cdot
\text{reverse}(s)$, $P(\varepsilon)$ holds.

\newpage
Inductive step: $t = y \cdot x$ where $x \in A$, $P(y) \to P(y \cdot x)$
\[
  P(y \cdot x) \DefEq \text{``for all $s \in A^*$ it holds
    that $\text{reverse}(s \cdot (y \cdot x)) = \text{reverse}(y \cdot x)
    \cdot \text{reverse}(s)$''}
\]

\begin{enumerate}
\item $\text{reverse}(s \cdot (y \cdot x)) = \text{reverse}(y \cdot x)
  \cdot \text{reverse}(s)$
\end{enumerate}

By the associative property of $\cdot$,

\begin{enumerate}[resume]
\item $\text{reverse}(s \cdot y \cdot x) = \text{reverse}(y \cdot x)
  \cdot \text{reverse}(s)$
\end{enumerate}

By $\text{reverse}(s \cdot x) \DefEq x \cdot \text{reverse}(s)$,

\begin{enumerate}[resume]
\item $x \cdot \text{reverse}(s \cdot y) = x \cdot \text{reverse}(y)
  \cdot \text{reverse}(s)$
\end{enumerate}

By the inductive hypothesis, $P(y)$,

\begin{enumerate}[resume]
\item $x \cdot \text{reverse}(y) \cdot \text{reverse}(s) = x \cdot
  \text{reverse}(y) \cdot \text{reverse}(s)$
\end{enumerate}

By the reflexive property, $P(y \cdot x)$ holds.

By structural induction, $P(t)$ holds. \hfill $\square$

\section*{Exercise 2}

[LCS, page 87]: Exercise 1.4.15

Use mathematical induction on $n$ to prove the theorem $(
  (\phi_1 \land (\phi_2 \land (\ldots \land \phi_n) \ldots) \to \psi)
    \to (\phi_1 \to (\phi_2 \to (\ldots (\phi_n \to \psi) \ldots)))
)$.

\medskip
Using natural deduction for $n = 2$ as a proof template:

\begin{logicproof}{3}
  \begin{subproof}
    (\phi_1 \land \phi_2) \to \psi & assumption \\
    \begin{subproof}
      \phi_1                       & assumption \\
      \begin{subproof}
        \phi_2                     & assumption \\
        \phi_1 \land \phi_2        & $\Intro{\land}$ 2, 3 \\
        \psi                       & $\Elim{\to}$ 1, 4
      \end{subproof}
      \phi_2 \to \psi              & $\Intro{\to}$ 3--5
    \end{subproof}
    \phi_1 \to (\phi_2 \to \psi)   & $\Intro{\to}$ 2--6
  \end{subproof}
  (\phi_1 \land \phi_2) \to \psi) \to (\phi_1 \to (\phi_2 \to \psi))
    & $\Intro{\to}$ 1--7
\end{logicproof}

\medskip
For the inductive proof we can recursively apply $\Intro{\to}$ until reaching
the premise $\phi_n$, then $\Intro{\land}$ $n - 1$ times to produce
$\phi_1 \land (\phi_2 \land (\ldots \land \phi_n) \ldots)$, then $\Elim{\to}$
to reach $\psi$.

\begin{logicproof}{3}
  \begin{subproof}
    (\phi_1 \land (\phi_2 \land (\ldots \land \phi_n) \ldots)) \to \psi
      & assumption \\
    \begin{subproof}
      \phi_1                               & assumption \\
      \vdots \vdots                        & \\
      \begin{subproof}
        \phi_n                             & assumption \\
        \phi_{n-1} \land \phi_n            & $\Intro{\land}$ 3, 4 \\
        \vdots \vdots                      & \\
        \phi_1 \land (\phi_2 \land (\ldots \land \phi_n) \ldots) & \\
        \psi                               & $\Elim{\to}$ 1, 7
      \end{subproof}
      \phi_n \to \psi                      & $\Intro{\to}$ 4--8 \\
      \vdots \vdots                        &
    \end{subproof}
    \phi_1 \to (\phi_2 \to (\ldots (\phi_n \to \psi) \ldots)) &
  \end{subproof}
  (\phi_1 \land (\phi_2 \land (\ldots \land \phi_n) \ldots) \to \psi)
    \to (\phi_1 \to (\phi_2 \to (\ldots (\phi_n \to \psi) \ldots)))
    & $\Intro{\to}$ 1--11
\end{logicproof}

\section*{Problem 1}

Show that any of the three rules $\{(LEM), (PBC), (\lnot\lnot E)\}$ are
interderivable.

% TODO

\vspace{3\baselineskip}
\hrule

Solutions to the Lean 4 exercises are available on GitHub:

\url{https://github.com/anthonyde/math2001/blob/cs511-fall2024/Math2001/CS511_Homework/hw2.lean}

\section*{Exercise 3}

\section*{Exercise 4}

\section*{Problem 2}

\end{document}
